% ============================================================
% Capítulo 2 — Referencial Teórico (versão reescrita)
% ParseH2IA — Guilherme Dallman Lima, 2026
% ============================================================

\chapter{Referencial Teórico}
\label{cap:referencial}

Este capítulo apresenta os fundamentos teóricos necessários para a compreensão do presente
trabalho. A organização segue a lógica da dissertação: primeiro, descreve-se a tarefa central
— a análise de dependências sintáticas e o padrão Universal Dependencies —; em seguida,
apresentam-se os modelos de linguagem pré-treinados baseados em Transformers e as
evidências de que esses modelos codificam estrutura sintática de forma linearmente acessível;
depois, discutem-se as arquiteturas de predição empregadas (linear e biaffine); por fim,
introduz-se o aprendizado multi-tarefa como paradigma que permite aproveitar a correlação entre
as tarefas auxiliar (UPOS) e principal (DEPREL/HEAD) investigadas neste trabalho.

% ============================================================
\section{Processamento de Linguagem Natural e Análise Sintática}
\label{sec:pln}
% ============================================================

O Processamento de Linguagem Natural (do inglês \textit{Natural Language Processing} —
NLP) é a área da inteligência artificial voltada ao tratamento e à interpretação automatizada de
texto e fala \cite{jurafsky2023speech}. Dentre as tarefas de NLP, a \textbf{análise sintática}
ocupa uma posição central: ela estrutura as relações gramaticais entre as palavras de uma
sentença e produz uma representação formal que serve de base para tarefas mais complexas,
como extração de informação, análise de sentimento baseada em aspecto, tradução automática e
resposta a perguntas. Neste trabalho, aborda-se especificamente a análise de dependências
sintáticas para o português brasileiro.

% ------------------------------------------------------------
\subsection{Análise de Constituência e Análise de Dependências}
\label{subsec:tipos-analise}
% ------------------------------------------------------------

A análise sintática pode ser conduzida segundo dois paradigmas principais, ilustrados na
Figura~\ref{fig:parsing_contituencia}.

Na \textbf{análise de constituência} (\textit{Constituency Parsing}), a sentença é segmentada
em sintagmas hierárquicos — Sintagma Nominal (SN), Sintagma Verbal (SV), etc. — que se
encaixam em uma árvore cujo nó raiz é a sentença inteira. A relação modelada é do tipo
parte-todo \cite{fernandez2022multitask}.

Na \textbf{análise de dependências} (\textit{Dependency Parsing}), a unidade central é a
palavra (\textit{token}), e as relações modeladas são binárias e direcionadas: cada palavra
dependente está ligada a exatamente uma palavra governante (\textit{head}), com uma etiqueta
que descreve a função sintática dessa relação (sujeito, objeto, modificador, etc.). O resultado
é uma árvore de dependências enraizada no verbo principal da sentença, conforme exemplificado
na Figura~\ref{fig:dep_tree_example}.

A análise de dependências tem três vantagens práticas para o presente trabalho: (i) representa
diretamente as relações gramaticais entre palavras, sem a necessidade de constituintes
intermediários; (ii) é mais robusta para línguas com ordem de palavras flexível, como o
português; e (iii) é o paradigma adotado pelo padrão Universal Dependencies, utilizado no
corpus Porttinari-base \cite{porttinari}.

\begin{figure}[htbp]
    \centering
    \includegraphics[width=1\textwidth]{Dissertação/Imagens/parser_constituency.png}
    \caption{Exemplo de árvore de análise de constituência para a sentença
             ``A menina brinca no parque.''}
    \label{fig:parsing_contituencia}
\end{figure}

% ------------------------------------------------------------
\subsection{Universal Dependencies e o Formato de Anotação}
\label{subsec:ud}
% ------------------------------------------------------------

O projeto \textit{Universal Dependencies} (UD) define um esquema de anotação morfossintática e
sintática com consistência interlinguística, atualmente abrangendo mais de cem línguas e
trezentos \textit{treebanks} \cite{UniversalDependencies}. O esquema se baseou nas dependências
universais de Stanford \cite{de2014universal}, nas etiquetas de classe gramatical do Google
\cite{petrov2011universal} e na Interlingua Interset para morfologia. O objetivo central é
disponibilizar categorias e diretrizes que permitam a anotação consistente de construções
semelhantes entre línguas distintas.

A Tabela~\ref{tab:ud-example} apresenta o formato CoNLL-U adotado pelo UD. Cada linha
representa um \textit{token} e contém dez campos: \textbf{ID} (posição na sentença),
\textbf{FORM} (forma superficial), \textbf{LEMMA} (forma canônica), \textbf{UPOS} (classe
morfossintática universal), \textbf{XPOS} (classe específica do idioma), \textbf{FEATS}
(atributos morfológicos, como gênero, número e tempo verbal), \textbf{HEAD} (índice do
governante), \textbf{DEPREL} (rótulo da relação de dependência), \textbf{DEPS} (dependências
secundárias) e \textbf{MISC} (informações adicionais).

\begin{table}[h]
\centering
\resizebox{1.05\textwidth}{!}{%
\renewcommand{\arraystretch}{2}
\sffamily
\begin{tabular}{
>{\centering\arraybackslash}p{0.2\textwidth}
>{\arraybackslash}p{0.15\textwidth}
>{\arraybackslash}p{0.1\textwidth}
>{\centering\arraybackslash}p{0.1\textwidth}
>{\centering\arraybackslash}p{0.1\textwidth}
>{\arraybackslash}p{0.1\textwidth}
>{\centering\arraybackslash}p{0.2\textwidth}
>{\arraybackslash}p{0.15\textwidth}
>{\arraybackslash}p{0.2\textwidth}
>{\arraybackslash}p{0.2\textwidth}
}
\toprule
\textbf{ID} & \textbf{FORM} & \textbf{LEMMA} & \textbf{UPOS} & \textbf{XPOS} &
\textbf{FEATS} & \textbf{HEAD} & \textbf{DEPREL} & \textbf{DEPS} & \textbf{MISC} \\
\midrule
1 & O & o & DET & \_ &
\makecell[l]{Definite=Def\\Gender=Masc\\Number=Sing\\PronType=Art} &
2 & det & 2:det & \_ \\
2 & canal & canal & NOUN & \_ &
\makecell[l]{Gender=Masc\\Number=Sing} &
3 & nsubj & 3:nsubj & \_ \\
3 & terá & ter & VERB & \_ &
\makecell[l]{Mood=Ind\\Number=Sing\\Person=3\\Tense=Fut\\VerbForm=Fin} &
0 & root & 0:root & \_ \\
4 & o & o & DET & \_ &
\makecell[l]{Definite=Def\\Gender=Masc\\Number=Sing\\PronType=Art} &
5 & det & 5:det & \_ \\
5 & conteúdo & conteúdo & NOUN & \_ &
\makecell[l]{Gender=Masc\\Number=Sing} &
3 & obj & \makecell[l]{3:obj\\6:nsubj:pass} & \_ \\
6 & reformulado & reformular & VERB & \_ &
\makecell[l]{Gender=Masc\\Number=Sing\\VerbForm=Part\\Voice=Pass} &
3 & xcomp & 3:xcomp & SpaceAfter=No \\
7 & . & . & PUNCT & \_ & \_ & 3 & punct & 3:punct & \_ \\
\bottomrule
\end{tabular}%
}
\caption{Sentença anotada no formato CoNLL-U (padrão Universal Dependencies).}
\label{tab:ud-example}
\end{table}

As relações entre palavras são representadas por conexões unidirecionadas, partindo de um
nó raiz (ROOT). Palavras que funcionam como núcleos (governadores) podem ter uma ou mais
palavras dependentes. A Figura~\ref{fig:dep_tree_example} ilustra a árvore de dependências
para a sentença ``A menina brinca no parque.'', com as etiquetas DEPREL anotadas nos arcos.

\begin{figure}[htbp]
    \centering
    \includegraphics[width=0.8\textwidth]{Dissertação/Imagens/arvore_dep.png}
    \caption{Árvore de dependências sintáticas para ``A menina brinca no parque.''
             As etiquetas nos arcos indicam as relações DEPREL: \texttt{det} (determinante),
             \texttt{nsubj} (sujeito nominal), \texttt{obl} (oblíquo), \texttt{case}
             (marcador de caso), \texttt{punct} (pontuação).}
    \label{fig:dep_tree_example}
\end{figure}

O projeto UD define 17 classes morfossintáticas universais (UPOS) e um conjunto de relações de
dependência (DEPREL) — como \texttt{nsubj}, \texttt{obj}, \texttt{nmod}, \texttt{obl},
\texttt{amod} — extensíveis por subtipos específicos de cada língua quando necessário. As
classes UPOS e os rótulos DEPREL utilizados neste trabalho seguem as diretrizes do
Porttinari-base \cite{porttinari}, detalhadas no Apêndice~\ref{ap:classes}.

% ------------------------------------------------------------
\subsection{Características do Português Brasileiro Relevantes para o Parsing}
\label{subsec:ptbr}
% ------------------------------------------------------------

O português brasileiro é uma língua morfologicamente rica: os verbos flexionam em pessoa,
número, tempo, modo, aspecto e voz; os nominais flexionam em gênero e número; e a posição
relativa dos argumentos é mais flexível do que em línguas de ordem rígida como o inglês.
Essa riqueza morfológica e a relativa liberdade de ordem de palavras trazem dois desafios
específicos para o parsing de dependências.

Primeiro, a desambiguação morfossintática é mais exigente: uma mesma forma superficial pode
corresponder a classes gramaticais distintas dependendo do contexto (por exemplo, ``para''
como verbo ou preposição), o que demanda representações contextualizadas para ser resolvida
com precisão. Segundo, relações de dependência de longa distância são mais comuns: fenômenos
como topicalização e deslocamento de constituintes criam dependências entre palavras distantes
na sequência linear, tornando mais difícil a predição do \textit{head} correto.

Estudos anteriores demonstraram que modelos monolíngues especializados no português, como o
BERTimbau \cite{souza2020bertimbau}, tendem a superar modelos multilíngues como o mBERT
\cite{devlin2019bert} em tarefas morfossintáticas para o português brasileiro, em parte por
capturarem padrões morfológicos específicos da língua durante o pré-treinamento em corpora
monolingues de maior qualidade \cite{lopes2024towards}.

% ============================================================
\section{Modelos de Linguagem Pré-treinados}
\label{sec:llm}
% ============================================================

O estado da arte em análise sintática para o português é dominado por abordagens baseadas em
\textit{encoders} Transformer pré-treinados. Esta seção descreve a arquitetura Transformer,
o modelo BERT e o processo de ajuste fino, que constituem a base dos quatro encoders avaliados
nesta dissertação.

% ------------------------------------------------------------
\subsection{Arquitetura Transformer}
\label{subsec:transformer}
% ------------------------------------------------------------

O Transformer \cite{vaswani2017attention} é uma arquitetura de redes neurais introduzida em
2017 para tarefas de tradução automática. Diferentemente de arquiteturas recorrentes (RNNs,
LSTMs), o Transformer processa toda a sequência de entrada em paralelo, capturando dependências
de longo alcance sem o problema do gradiente evanescente. A arquitetura completa, composta por
um codificador e um decodificador, é ilustrada na Figura~\ref{fig:transformer}.

\begin{figure}[htbp]
    \centering
    \includegraphics[width=0.65\textwidth]{Dissertação/Imagens/transformer.png}
    \caption{Arquitetura Transformer \cite{vaswani2017attention}. O codificador (esquerda)
             processa a sequência de entrada; o decodificador (direita) gera a sequência de
             saída condicionada ao codificador.}
    \label{fig:transformer}
\end{figure}

O bloco fundamental do Transformer é a \textbf{atenção multi-cabeça}
(\textit{multi-head attention}). Dado um conjunto de consultas $Q$, chaves $K$ e valores $V$,
a atenção escalada por produto escalar é definida como:

\begin{equation}
\text{Attention}(Q,K,V) = \text{softmax}\!\left(\frac{QK^\top}{\sqrt{d_k}}\right)V,
\end{equation}

\noindent
onde $d_k$ é a dimensão dos vetores de chave e o fator $1/\sqrt{d_k}$ evita que os produtos
escalares cresçam para valores que saturem a softmax. A atenção multi-cabeça executa $h$
atenções em paralelo sobre subespaços projetados e concatena os resultados:

\begin{equation}
\text{MultiHead}(Q,K,V) = \text{Concat}(\text{head}_1,\ldots,\text{head}_h)\,W_O,
\end{equation}

\noindent
onde cada $\text{head}_i = \text{Attention}(QW_i^Q,\, KW_i^K,\, VW_i^V)$. Isso permite que
o modelo atenda simultaneamente a informações em diferentes posições e subespaços de
representação.

Como o Transformer não utiliza recorrência, a ordem da sequência é injetada por meio de
\textbf{codificação posicional}, somada aos \textit{embeddings} de entrada:

\begin{equation}
PE_{(pos,\,2i)}   = \sin\!\left(\frac{pos}{10000^{2i/d_{\text{model}}}}\right),\quad
PE_{(pos,\,2i+1)} = \cos\!\left(\frac{pos}{10000^{2i/d_{\text{model}}}}\right).
\end{equation}

Cada camada do codificador contém, além da atenção multi-cabeça, uma rede \textit{feed-forward}
totalmente conectada aplicada posição a posição, seguida de normalização de camada e conexão
residual. As saídas de cada camada mantêm a mesma dimensionalidade, permitindo o empilhamento
de $N$ blocos idênticos para maior capacidade de abstração.

% ------------------------------------------------------------
\subsection{BERT: Pré-treino Bidirecional}
\label{subsec:bert}
% ------------------------------------------------------------

O BERT (\textit{Bidirectional Encoder Representations from Transformers}) \cite{devlin2019bert}
é um \textit{encoder} Transformer pré-treinado que utiliza exclusivamente o componente de
codificação, sem decodificador. Isso o torna apropriado para tarefas de compreensão de
linguagem, onde o objetivo é produzir representações ricas para cada token da entrada, e não
gerar sequências novas.

O BERT é pré-treinado por dois objetivos complementares:

\begin{itemize}
    \item \textbf{Modelagem de Linguagem Mascarada} (\textit{Masked Language Modeling} —
    MLM): aproximadamente 15\% dos tokens de entrada são substituídos por um token especial
    \texttt{[MASK]}, e o modelo é treinado a prever os tokens originais a partir do contexto
    bidirecional restante. Ao contrário de modelos autorregressivos (da esquerda para a
    direita), o MLM permite que cada posição utilize informação tanto do lado esquerdo quanto
    do direito, produzindo representações bidirecionalmente contextualizadas.

    \item \textbf{Predição da Próxima Sentença} (\textit{Next Sentence Prediction} — NSP):
    o modelo aprende a classificar se duas sentenças ocorrem consecutivamente no corpus
    original, adquirindo noções de coerência discursiva.
\end{itemize}

O BERT opera com contexto de até 512 tokens. Tokens especiais \texttt{[CLS]} e \texttt{[SEP]}
delimitam as entradas. A saída do encoder é um tensor $\mathbf{H} \in \mathbb{R}^{B \times L
\times d}$, onde $B$ é o tamanho do lote, $L$ o comprimento máximo da sequência e $d$ a
dimensão oculta ($d = 768$ para modelos Base e $d = 1024$ para modelos Large). As saídas
posicionais $\mathbf{h}_i$ são utilizadas diretamente como representações dos tokens nas
tarefas de anotação.

Neste trabalho são avaliados quatro \textit{encoders}: mBERT (multilíngue)
\cite{devlin2019bert}, BERTimbau Base e BERTimbau Large (monolíngues para o português
brasileiro) \cite{souza2020bertimbau}, e JabuticaBERT (monolíngue, com vocabulário expandido
baseado no ModernBERT) \cite{thiago}.

% ------------------------------------------------------------
\subsection{Ajuste Fino}
\label{subsec:finetuning}
% ------------------------------------------------------------

O \textbf{ajuste fino} (\textit{fine-tuning}) consiste em inicializar um modelo com os pesos
pré-treinados e continuar o treinamento sobre dados anotados para a tarefa alvo
\cite{devlin2019bert}. Durante esse processo, todos os parâmetros — tanto os da tarefa quanto
os do encoder — são atualizados conjuntamente por retropropagação. O ajuste fino completo
(\textit{full fine-tuning}) é o regime adotado nesta dissertação: o encoder não é mantido
congelado, e seus parâmetros são adaptados junto com os cabeçotes de predição. Essa
abordagem contrasta com estratégias de adaptação mais parcimoniosas (como \textit{adapters}
ou \textit{LoRA}), mas é justificada pelo corpus de dimensão moderada e pela necessidade de
máxima especialização das representações para tarefas sintáticas.

% ------------------------------------------------------------
\subsection{O que o BERT Aprende: Evidências de Probing Sintático}
\label{subsec:probing}
% ------------------------------------------------------------

Uma questão central para esta dissertação é: \emph{em que medida as representações produzidas
por um encoder BERT já codificam estrutura sintática, e em que forma essa informação está
disponível?} Essa pergunta tem motivação direta nos resultados da dissertação — a
\textbf{superioridade consistente do cabeçote linear sobre o biaffine} — e encontra respaldo
em uma série de estudos de \textit{probing} (investigação de representações).

\textbf{Tenney et al.} \cite{tenney2019bert} investigaram em quais camadas do BERT diferentes
tipos de informação linguística estão codificados. Usando classificadores lineares (\textit{edge
probes}) treinados sobre as representações de cada camada, os autores mostraram que o BERT
``redesobre o pipeline clássico de NLP'': informação morfossintática (POS) emerge nas camadas
iniciais, informação sintática local (dependências de curto alcance, constituência) concentra-se
nas camadas intermediárias, e informação semântica de ordem superior aparece nas camadas mais
profundas. Crucialmente, a performance dos probes lineares nas camadas intermediárias para
tarefas sintáticas é alta, indicando que essa informação está linearmente acessível nas
representações.

\textbf{Hewitt e Manning} \cite{hewitt2019structural} propuseram uma sonda estrutural
(\textit{structural probe}) para verificar se as distâncias entre pares de palavras no espaço
de representação do BERT são proporcionais às distâncias na árvore de dependências sintáticas.
A hipótese é que, se existir uma transformação linear $B$ tal que
$\lVert B(\mathbf{h}_i - \mathbf{h}_j) \rVert_2 \approx d_{\text{tree}}(i, j)$, então a
estrutura sintática está geometricamente codificada de forma linearmente acessível. Os
resultados mostraram que tal transformação existe e recupera com alta precisão a árvore de
dependências a partir das representações do BERT, sem treinamento adicional além da sonda
linear. Isso fornece evidência direta de que \textbf{a estrutura de dependências está
implicitamente representada no espaço linear do encoder}.

\textbf{Jawahar et al.} \cite{jawahar2019does} confirmaram resultados complementares: as
camadas intermediárias do BERT capturam informações sobre estrutura frasal (sintagmas) e
hierarquia sintática, enquanto as camadas superiores processam relações semânticas mais
abstratas. Probes lineares treinados sobre representações extraídas diretamente do encoder
classificam relações de dependência com acurácia substancialmente acima do acaso.

Em conjunto, esses trabalhos embasam a hipótese central desta dissertação: \textbf{o encoder
BERT, após ajuste fino completo, organiza suas representações de modo que a estrutura
sintática seja linearmente acessível, favorecendo cabeçotes de predição simples (lineares)
como sondas diretas dessa organização latente}. Isso explica por que o cabeçote linear, apesar
de parametricamente mais simples que o biaffine, consegue superar este último quando o corpus
de treinamento é limitado: ele extrai eficientemente a informação sintática já disponível no
encoder, enquanto o cabeçote biaffine — com sua maior capacidade paramétrica — incorre em
maior risco de sobreajuste.

% ============================================================
\section{Arquiteturas de Predição para Parsing de Dependências}
\label{sec:cabecotes}
% ============================================================

O parsing de dependências envolve duas predições por token: (i) qual token é o governante
(\textit{head}) — tarefa HEAD — e (ii) qual é a relação sintática entre o token e seu
governante — tarefa DEPREL. Esta seção descreve as duas arquiteturas de cabeçote avaliadas
neste trabalho: o cabeçote linear e o cabeçote biaffine.

% ------------------------------------------------------------
\subsection{Cabeçote Linear}
\label{subsec:linear}
% ------------------------------------------------------------

No cabeçote linear, a saída do encoder $\mathbf{H} \in \mathbb{R}^{B \times L \times d}$ é
projetada por camadas densas independentes para cada tarefa. Para a tarefa HEAD, a projeção
mapeia cada representação de token para uma distribuição sobre $N_{\text{head}}$ posições
candidatas; para DEPREL, para os 44 rótulos de dependência do Porttinari; e para UPOS, para
os $N_u$ rótulos morfossintáticos. A predição de HEAD no cabeçote linear trata a tarefa como
uma classificação sobre posições absolutas da sentença, sem modelar explicitamente a interação
entre pares de tokens.

Essa arquitetura tem capacidade paramétrica reduzida, o que atua como regularização implícita
em cenários de dados limitados. Ao mesmo tempo, conforme discutido na
Seção~\ref{subsec:probing}, as representações do BERT já codificam informação sintática em
forma linearmente acessível, de modo que uma projeção linear pode ser suficiente para extrair
essa estrutura latente sem necessidade de aprendizado adicional de interações entre pares.

% ------------------------------------------------------------
\subsection{Atenção Biaffine para Parsing de Dependências}
\label{subsec:biaffine}
% ------------------------------------------------------------

O mecanismo de atenção biaffine para parsing de dependências foi proposto por
\citeauthor{dozat2016deep} \cite{dozat2016deep} e representa o estado da arte em sistemas de
parsing graph-based. A ideia central é modelar explicitamente a compatibilidade entre cada
par de tokens na sentença, projetando as representações do encoder para papéis especializados
antes de calcular os escores de dependência.

Para cada token, quatro MLPs especializadas projetam a representação do encoder para papéis
distintos: cabeça de arco, dependente de arco, cabeça de relação e dependente de relação.
Em seguida, um operador bilinear calcula, para cada par de tokens $(i,j)$, um escore que
modela a compatibilidade entre a representação do dependente $i$ e a do candidato a cabeça
$j$ — tanto para a existência do arco quanto para o rótulo da relação. A implementação
detalhada, incluindo as equações e as dimensões das matrizes de peso, é apresentada na
Seção~\ref{subsec:biaffine_met}.

A operação bilinear modela interações multiplicativas entre as representações do dependente
e da cabeça, o que torna o mecanismo estritamente mais expressivo do que uma projeção linear
sobre a representação de um único token. Em corpus de grande porte, essa expressividade
adicional resulta em ganhos consistentes de desempenho. Entretanto, a maior capacidade
paramétrica aumenta o risco de sobreajuste em corpus pequenos, como o Porttinari-base
(5.893 sentenças de treinamento), motivando a comparação direta com o cabeçote linear
conduzida nesta dissertação.

% ============================================================
\section{Aprendizado Multi-Tarefa}
\label{sec:mtl}
% ============================================================

O \textbf{aprendizado multi-tarefa} (\textit{Multi-Task Learning} — MTL) é um paradigma que
treina um único modelo para resolver simultaneamente múltiplas tarefas relacionadas,
aproveitando a correlação entre elas para melhorar a generalização em cada tarefa individual
\cite{caruana1997multitask}. Em vez de treinar modelos independentes para cada tarefa, o MTL
induz ao compartilhamento de representações internas, funcionando como uma forma de
regularização implícita que reduz o risco de sobreajuste.

% ------------------------------------------------------------
\subsection{Hard Parameter Sharing}
\label{subsec:hard-sharing}
% ------------------------------------------------------------

O \textit{hard parameter sharing} é a estratégia de MTL adotada nesta dissertação. Nela, um
único encoder compartilha todos os seus parâmetros entre as tarefas, enquanto cada tarefa
possui seu próprio cabeçote de predição. A função de perda total é a soma das perdas
individuais:

\begin{equation}
\mathcal{L} = \mathcal{L}_{\text{HEAD}} + \mathcal{L}_{\text{DEPREL}} + \mathcal{L}_{\text{UPOS}}.
\end{equation}

\noindent
Dessa forma, os gradientes de todas as tarefas propagam-se pelos mesmos parâmetros do encoder,
forçando as representações a capturarem características úteis para cada uma delas. Quando as
tarefas estão correlacionadas, esse compartilhamento melhora as representações de cada tarefa
individual; quando as tarefas são independentes ou conflitantes, o compartilhamento pode
prejudicar o desempenho.

A alternativa — o \textit{soft parameter sharing} — mantém encoders separados por tarefa e
introduz um mecanismo de regularização para aproximar seus parâmetros. Embora mais flexível,
essa abordagem é parametricamente mais custosa e não foi adotada neste trabalho.

% ------------------------------------------------------------
\subsection{UPOS como Tarefa Auxiliar: Motivação para o MTL}
\label{subsec:mtl-upos}
% ------------------------------------------------------------

A escolha de UPOS como tarefa auxiliar para o parsing de dependências é motivada pela
correlação linguística entre classe morfossintática e função sintática. Intuitivamente,
saber que uma palavra é um verbo favorece a predição de que ela seja raiz ou núcleo verbal,
enquanto saber que é um pronome pessoal favorece a predição de sujeito ou objeto. Essa
correlação torna as tarefas estatisticamente interdependentes: a distribuição de rótulos
DEPREL condicionada ao UPOS não é uniforme.

A dependência estatística entre UPOS e DEPREL é verificada empiricamente no corpus
Porttinari-base por meio de métricas como Informação Mútua Normalizada (NMI) e V de Cramér,
cuja formalização e resultados numéricos são apresentados na Seção~\ref{subsec:corr}. Os
resultados confirmam forte associação entre as duas variáveis e fundamentam quantitativamente
o uso de UPOS como sinal auxiliar no treinamento conjunto.

Trabalhos anteriores em línguas morfologicamente ricas confirmam que informações
morfossintáticas compartilhadas em um modelo MTL melhoram o desempenho de parsing de
dependências, especialmente para rótulos de menor frequência
\cite{kondratyuk201975}. Nesta dissertação, adicionalmente investiga-se o impacto da
remoção do sinal UPOS por meio de um estudo de ablação, detalhado na
Seção~\ref{subsec:ablacao}.
